% ==============================================================================
% CONFIGURACIÓN Y ESTILO VISUAL — FUNDAMENTOS FÍSICOS Y TECNOLÓGICOS (FFT)
% Escuela Técnica Superior de Ingenierías Informática y de Telecomunicación
% Universidad de Granada (UGR) — Plataforma repasaYA
% ==============================================================================

\usepackage[utf8]{inputenc}
\usepackage[T1]{fontenc}
\usepackage[spanish,es-tabla]{babel}
\usepackage{lmodern}
\usepackage{microtype}
\usepackage[dvipsnames,svgnames,table,xcdraw]{xcolor}
\usepackage{amsmath,amssymb,amsfonts,amsthm}
\usepackage{physics}
\usepackage{siunitx}
\usepackage{graphicx}
\usepackage{float}
\usepackage{booktabs}
\usepackage{array}
\usepackage{multirow}
\usepackage{enumitem}
\usepackage{fancyhdr}
\usepackage{titlesec}
\usepackage{listings}
\usepackage{varwidth}
\usepackage{tcolorbox}
\tcbuselibrary{skins,breakable}
\usepackage{tikz}
\usetikzlibrary{babel,arrows.meta,positioning,calc,patterns}
\usepackage[siunitx,european,cuteinductors]{circuitikz}

% Configuración de márgenes
\usepackage[a4paper,top=2.5cm,bottom=2.5cm,left=2.8cm,right=2.8cm,headheight=24pt]{geometry}

% Configuración de siunitx en español
\sisetup{
  output-decimal-marker = {,},
  range-phrase = { a },
  per-mode = symbol
}
\DeclareSIUnit\VAR{VAr}
\DeclareSIUnit\VA{VA}

% Paleta de colores profesionales para Física y Electrónica
\definecolor{fftDarkBlue}{RGB}{16, 32, 60}
\definecolor{fftPrimary}{RGB}{24, 76, 120}
\definecolor{fftAccent}{RGB}{210, 85, 20}
\definecolor{fftLightAccent}{RGB}{254, 245, 238}
\definecolor{fftGreen}{RGB}{28, 130, 65}
\definecolor{fftLightGreen}{RGB}{240, 253, 244}
\definecolor{fftOrange}{RGB}{210, 110, 10}
\definecolor{fftLightOrange}{RGB}{255, 251, 235}
\definecolor{fftRed}{RGB}{185, 28, 28}
\definecolor{fftLightRed}{RGB}{254, 242, 242}
\definecolor{fftGray}{RGB}{246, 248, 250}
\definecolor{fftBorderGray}{RGB}{215, 222, 232}
\definecolor{fftTeal}{RGB}{15, 115, 110}
\definecolor{codeBg}{RGB}{248, 249, 252}

% Hipervínculos
\usepackage{hyperref}
\hypersetup{
    colorlinks=true,
    linkcolor=fftDarkBlue,
    citecolor=fftTeal,
    urlcolor=fftAccent,
    pdfauthor={repasaYA - UGR},
    pdftitle={Fundamentos Físicos y Tecnológicos - Apuntes Completos},
    pdfsubject={Electromagnetismo, Circuitos DC/AC, Semiconductores, Diodos, MOSFET y Operacionales},
    pdfkeywords={FFT, Electromagnetismo, Circuitos, Kirchhoff, Thévenin, Diodos, MOSFET, CMOS, Op-Amp, UGR}
}

% Configuración de Encabezados y Pies de página
\pagestyle{fancy}
\fancyhf{}
\fancyhead[LE,RO]{\footnotesize\textsf{\thepage}}
\fancyhead[LO]{\footnotesize\textsf{\rightmark}}
\fancyhead[RE]{\footnotesize\textsf{\leftmark}}
\renewcommand{\headrulewidth}{0.4pt}
\renewcommand{\footrulewidth}{0pt}

% Configuración de títulos de capítulos y secciones
\titleformat{\chapter}[display]
  {\normalfont\bfseries\color{fftDarkBlue}\sffamily}
  {\Large\textsf{Tema \thechapter}}{1ex}
  {\titlerule[1.5pt]\vspace{1ex}\huge}
  [\vspace{1ex}\titlerule]

\titleformat{\section}
  {\normalfont\Large\bfseries\color{fftPrimary}\sffamily}
  {\thesection}{1em}{}

\titleformat{\subsection}
  {\normalfont\large\bfseries\color{fftDarkBlue}\sffamily}
  {\thesubsection}{1em}{}

\titleformat{\subsubsection}
  {\normalfont\normalsize\bfseries\color{fftAccent}\sffamily}
  {\thesubsubsection}{1em}{}

% Entorno: Caja de Fuentes y Metadatos al inicio de cada tema
\newcommand{\cuadrofuentes}[1]{%
  \begin{tcolorbox}[
    colback=fftLightAccent,
    colframe=fftAccent,
    boxrule=1pt,
    arc=4pt,
    left=12pt,
    right=12pt,
    top=8pt,
    bottom=8pt,
    breakable
  ]
    \small
    \noindent\textbf{\color{fftAccent}\large\textsf{Metadatos del Tema y Fuentes:}}\\[4pt]
    #1
  \end{tcolorbox}
  \vspace{10pt}
}

% Entorno: Definición
\newtcolorbox{definicion}[1][]{
  enhanced,
  colback=fftGray,
  colframe=fftPrimary,
  boxrule=1pt,
  leftrule=4pt,
  arc=3pt,
  title=\textbf{\textsf{Definición\ifstrempty{#1}{}{: #1}}},
  coltitle=white,
  colbacktitle=fftPrimary,
  attach boxed title to top left={yshift=-2mm, xshift=4mm},
  boxed title style={arc=2pt, boxrule=0pt},
  varwidth boxed title*=-8mm,
  top=10pt,
  bottom=8pt,
  left=10pt,
  right=10pt,
  breakable
}

% Entorno: Ley / Teorema Fundamental
\newtcolorbox{ley}[1][]{
  enhanced,
  colback=fftLightAccent,
  colframe=fftAccent,
  boxrule=0.8pt,
  leftrule=3.5pt,
  arc=3pt,
  title=\textbf{\textsf{Ley / Teorema\ifstrempty{#1}{}{: #1}}},
  coltitle=white,
  colbacktitle=fftAccent,
  attach boxed title to top left={yshift=-2mm, xshift=4mm},
  boxed title style={arc=2pt, boxrule=0pt},
  varwidth boxed title*=-8mm,
  top=10pt,
  bottom=8pt,
  left=10pt,
  right=10pt,
  breakable
}

% Entorno: Ejemplo Explicado
\newtcolorbox{ejemplo}[1][]{
  enhanced,
  colback=fftLightGreen,
  colframe=fftGreen,
  boxrule=0.8pt,
  leftrule=3.5pt,
  arc=3pt,
  title=\textbf{\textsf{Ejemplo Explicado\ifstrempty{#1}{}{: #1}}},
  coltitle=white,
  colbacktitle=fftGreen,
  attach boxed title to top left={yshift=-2mm, xshift=4mm},
  boxed title style={arc=2pt, boxrule=0pt},
  varwidth boxed title*=-8mm,
  top=10pt,
  bottom=8pt,
  left=10pt,
  right=10pt,
  breakable
}

% Entorno: Aviso / Ojo Examen
\newtcolorbox{aviso}[1][]{
  enhanced,
  colback=fftLightOrange,
  colframe=fftOrange,
  boxrule=0.8pt,
  leftrule=3.5pt,
  arc=3pt,
  title=\textbf{\textsf{Importante / Ojo Examen\ifstrempty{#1}{}{: #1}}},
  coltitle=white,
  colbacktitle=fftOrange,
  attach boxed title to top left={yshift=-2mm, xshift=4mm},
  boxed title style={arc=2pt, boxrule=0pt},
  varwidth boxed title*=-8mm,
  top=10pt,
  bottom=8pt,
  left=10pt,
  right=10pt,
  breakable
}

% Entorno: Método Sistemático de Resolución
\newtcolorbox{metodo}[1][]{
  enhanced,
  colback=codeBg,
  colframe=fftDarkBlue,
  boxrule=0.8pt,
  leftrule=3.5pt,
  arc=3pt,
  title=\textbf{\textsf{Método Sistemático\ifstrempty{#1}{}{: #1}}},
  coltitle=white,
  colbacktitle=fftDarkBlue,
  attach boxed title to top left={yshift=-2mm, xshift=4mm},
  boxed title style={arc=2pt, boxrule=0pt},
  varwidth boxed title*=-8mm,
  top=10pt,
  bottom=8pt,
  left=10pt,
  right=10pt,
  breakable
}

% Entorno: Ejercicio Resuelto
\newtcolorbox{ejercicio}[1][]{
  enhanced,
  colback=fftGray,
  colframe=fftTeal,
  boxrule=0.8pt,
  leftrule=3.5pt,
  arc=3pt,
  title=\textbf{\textsf{Ejercicio Resuelto\ifstrempty{#1}{}{: #1}}},
  coltitle=white,
  colbacktitle=fftTeal,
  attach boxed title to top left={yshift=-2mm, xshift=4mm},
  boxed title style={arc=2pt, boxrule=0pt},
  varwidth boxed title*=-8mm,
  top=10pt,
  bottom=8pt,
  left=10pt,
  right=10pt,
  breakable
}

% Configuración de listings
\lstset{
  backgroundcolor=\color{codeBg},
  basicstyle=\ttfamily\footnotesize\color{black},
  keywordstyle=\bfseries\color{fftPrimary},
  commentstyle=\itshape\color{gray},
  stringstyle=\color{fftGreen},
  numbers=left,
  numberstyle=\tiny\color{gray},
  numbersep=8pt,
  stepnumber=1,
  breaklines=true,
  breakatwhitespace=false,
  showstringspaces=false,
  frame=none,
  tabsize=2,
  captionpos=b,
  inputencoding=utf8,
  extendedchars=true,
  literate=
    {á}{{\'a}}1 {é}{{\'e}}1 {í}{{\'i}}1 {ó}{{\'o}}1 {ú}{{\'u}}1
    {Á}{{\'A}}1 {É}{{\'E}}1 {Í}{{\'I}}1 {Ó}{{\'O}}1 {Ú}{{\'U}}1
    {ñ}{{\~n}}1 {Ñ}{{\~N}}1 {¿}{{?`}}1 {¡}{{!`}}1
}
